% ---
% title: Odds in betting and prediction markets
% tags: finance, odds
% ---

\begin{problem}
	\begin{context}
		On a \emph{prediction market}, people trade contracts that pay 1 euro if
		an event happens (for example, a candidate wins an election) and nothing
		otherwise. The price of such a contract is read as the market's
		probability of the event. Bookmakers quote \emph{decimal odds}: the total
		amount a stake of 1 euro returns if the bet wins, the stake included.
	\end{context}
	(a) A prediction-market contract pays 1 euro if a candidate wins the
	election and 0 otherwise. It trades at 0.40 euros. What probability and
	what odds does the market assign to her winning?
	(b) A bookmaker offers decimal odds 2.50 on team A and 1.60 on team B in a
	match that cannot end in a draw: a stake of 1 euro returns 2.50 euros on
	team A, or 1.60 euros on team B, stake included, if the bet wins. What
	probabilities
	do these quotes imply? Why do they not add up to 1?
	(c) News arrives that is 3 times as likely if the candidate in (a) wins as
	if she loses. What is the fair price of the contract now?
\end{problem}

\begin{slideproblem}
	(a) A contract pays 1 euro if a candidate wins the election, 0 otherwise;
	its price is 0.40 euros. What probability and what odds does the market
	give to her winning?
	(b) A bookmaker quotes decimal odds 2.50 on team A and 1.60 on team B (no
	draw; a 1-euro stake returns 2.50 euros on A, or 1.60 euros on B, if it
	wins). Implied
	probabilities? Why is their sum not 1?
	(c) News arrives that is 3 times as likely if the candidate wins as if she
	loses. New fair price of the contract?
\end{slideproblem}

\begin{slides}
	\small
	\exstatement
	\pause
	\sol (a) $W$ = ``she wins''. A contract that pays 1 euro if $W$ happens
	is priced at the probability of $W$: \ $P(W)=0.40$.\\[1mm]
	\pause
	Odds of $W$: \
	$\dfrac{P(W)}{P(W^c)}=\dfrac{0.40}{0.60}=\dfrac23$, \ ``2 to 3 in favor''
	(3 to 2 against).\\[1mm]
	\pause
	(b) Decimal odds $d$: a stake of $1/d$ euros returns exactly 1 euro. So
	$1/d$ is the price of ``1 euro if the team wins'', the implied
	probability.\\[1mm]
	\pause
	Team A: $\dfrac{1}{2.50}=0.40$, \qquad team B: $\dfrac{1}{1.60}=0.625$.
\end{slides}

\begin{slides}
	\small
	(b) Implied probabilities $0.40$ and $0.625$.\\[1mm]
	\pause
	Their sum is $1.025>1$. Bet $0.40$ on A and $0.625$ on B: you get 1 euro
	back for sure, but you paid $1.025$ euros. The extra $0.025$ (2.5\%) is
	the bookmaker's margin (the \emph{overround}).\\[1mm]
	\pause
	Remove the margin by dividing by the sum:
	\[ P(\text{A wins})\approx\frac{0.40}{1.025}\approx 0.390,\qquad
		P(\text{B wins})\approx\frac{0.625}{1.025}\approx 0.610. \]
	\pause
	(c) Likelihood ratio of the news: $3$. New odds $=$ old odds $\times$
	likelihood ratio:
	\[ \frac23\cdot 3=2. \]
	\pause
	Back to a probability:
	$P(W\mid\text{news})=\dfrac{2}{1+2}=\dfrac23\approx 0.667$. \ The fair
	price rises from $0.40$ to about $0.667$ euros.
\end{slides}

\begin{solution}
	\idea A contract that pays 1 euro if an event happens is worth the
	probability of the event, so prices and decimal odds can be turned into
	probabilities. Then we use the odds of an event,
	\[ \text{odds of } H=\frac{P(H)}{P(H^c)},\qquad
		P(H)=\frac{\text{odds}}{1+\text{odds}}, \]
	and Bayes' formula in odds form: new odds $=$ old odds $\times$ likelihood
	ratio.

	\step{(a)} Let $W$ be the event ``the candidate wins''. The contract pays 1
	euro if $W$ happens and 0 otherwise. Suppose the true probability of $W$ is
	$q$. If the price were below $q$, buying contracts would be a good deal on
	average; if it were above $q$, selling them would be. Trading pushes the
	price towards $q$, so the price is read as the market's probability:
	\[ P(W)=0.40. \]
	The odds of $W$ compare the probability that she wins with the probability
	that she loses:
	\[ \frac{P(W)}{P(W^c)}=\frac{0.40}{1-0.40}=\frac{0.40}{0.60}=\frac23. \]
	In words: ``2 to 3 in favor'', or equivalently ``3 to 2 against''. Out of
	every 5 equally likely scenarios, she wins in 2 and loses in 3.

	\step{(b)} With decimal odds $d$, a stake of 1 euro returns $d$ euros if
	the bet wins. So a stake of $1/d$ euros returns exactly 1 euro if the bet
	wins, and nothing otherwise. This is the contract of part (a), and its
	price $1/d$ is the implied probability:
	\[ \text{team A: } \frac{1}{2.50}=0.40,\qquad
		\text{team B: } \frac{1}{1.60}=0.625. \]
	The match has no draw, so exactly one team wins: the events ``A wins'' and
	``B wins'' are mutually exclusive and together cover every outcome. Their
	true probabilities must therefore add up to 1. The quotes, however, give
	\[ 0.40+0.625=1.025>1. \]
	The reason is the bookmaker's margin. Bet 0.40 euros on A and 0.625 euros
	on B. Whichever team wins, you get exactly 1 euro back, but you paid
	$1.025$ euros: you lose $0.025$ euros for sure. This extra 2.5\% is the
	bookmaker's profit, called the \emph{overround}. A fair bookmaker would set
	the quotes so that the implied probabilities add up to exactly 1.

	To estimate the probabilities behind the quotes, remove the margin by
	dividing each implied probability by their sum (then they add up to 1):
	\[ P(\text{A wins})\approx\frac{0.40}{1.025}\approx 0.390,\qquad
		P(\text{B wins})\approx\frac{0.625}{1.025}\approx 0.610. \]

	\step{(c)} Call the news $N$. We are told that
	$P(N\mid W)=3\,P(N\mid W^c)$, so the likelihood ratio of the news is
	\[ \frac{P(N\mid W)}{P(N\mid W^c)}=3. \]
	We do not need the two conditional probabilities themselves, only their
	ratio. Bayes' formula in odds form gives
	\[ \frac{P(W\mid N)}{P(W^c\mid N)}=\frac{P(W)}{P(W^c)}\cdot
		\frac{P(N\mid W)}{P(N\mid W^c)}=\frac23\cdot 3=2. \]
	The new odds are 2, ``2 to 1 in favor''. Back to a probability:
	\[ P(W\mid N)=\frac{2}{1+2}=\frac23\approx 0.667. \]
	The contract still pays 1 euro if she wins, so its fair price is now about
	$0.667$ euros (up from $0.40$ euros).
	\begin{answer}
		(a) probability $0.40$, odds $\frac23$ (2 to 3 in favor) \quad
		(b) $0.40$ and $0.625$; the sum $1.025$ includes a 2.5\% margin;
		without it about $0.390$ and $0.610$ \quad
		(c) $\frac23\approx 0.667$ euros
	\end{answer}
	\begin{remark}
		\textbf{Common mistake:} multiplying the \emph{probability} by the
		likelihood ratio: $0.40\cdot 3=1.2$, which is not even a probability.
		The likelihood ratio multiplies the odds; only at the end do we turn the
		odds back into a probability.
	\end{remark}
	\begin{remark}
		The ``odds'' of a bookmaker are not the odds of this course. Decimal odds
		2.50 mean a net gain of 1.50 euros per euro staked, that is ``1.5 to 1
		against'' (3 to 2 against). In our sense the odds of team A winning are
		$0.40/0.60=2/3$, exactly as for the candidate in (a).
	\end{remark}
\end{solution}
